% GENERATED FILE. Edit canonical lessons, then run npm run generate:books.
% canonical-source-hash: fnv1a64:372e3f940d9c5f07
% canonical-lessons: PT-C05-falar, PT-C05-morar, PT-C05-trabalhar, PT-C05-falo-portugues, PT-C05-practice

\chapter{The First Verbs}
\label{ch:pt-first-verbs}
\hlchaptermodality{5}

\begin{chapteropening}
\textbf{By the end of this chapter:} \emph{I can build my own Portuguese sentences with -ar verbs and say what I speak, where I live and where I work.}

Conjugate falar, morar and trabalhar across eu, tu and você with the pronoun dropped, then answer Você fala português? and Onde você mora? about yourself.
\end{chapteropening}

\section[falar]{falar — \textquotedblleft{}to speak,\textquotedblright{} and the \emph{f} that Spanish lost}

\label{lesson:PT-C05-falar}



\begin{quote}
Until now you've assembled fixed phrases. \textbf{Falar} (“to speak”) is your first real \textbf{verb} and the model for the large \textbf{-ar} family. Its first letter also hides a beautiful contrast with Spanish.
\end{quote}

\begin{sounds}
\begin{itemize}
\raggedright
  \item \texttt{final-r} — \textbf{falar} = \emph{fah-LAR}, tapped final \emph{r}. The stem is \emph{fal-}.
\end{itemize}
\end{sounds}

\begin{cousinweb}
\textbf{falar} comes from Latin \textbf{fabulārī}, \emph{\textquotedblleft{}to chat, to tell tales\textquotedblright{}} — from \textbf{fābula}, \textquotedblleft{}a story\textquotedblright{} ($\to$ English \textbf{fable}, \textbf{fabulous}, \textbf{confabulate}). The \textbf{same Latin verb as Spanish \emph{hablar}} — but here's the twist:

\textbf{Portuguese kept the Latin \emph{f-}; Spanish turned \emph{f-} into \emph{h-}.} So the two \textquotedblleft{}speak\textquotedblright{} verbs diverge at letter one, and the pattern runs through the whole vocabulary:

\noindent
\begin{tabularx}{\linewidth}{@{}>{\raggedright\arraybackslash}X>{\raggedright\arraybackslash}X>{\raggedright\arraybackslash}X@{}}
\toprule
\textbf{Latin (f-)} & \textbf{Portuguese (f-)} & \textbf{Spanish (h-, silent)} \\
\midrule
\emph{fabulārī} & \textbf{falar} & hablar \\
\emph{farīna} & \textbf{farinha} (flour) & harina \\
\emph{fīlius} & \textbf{filho} (son) & hijo \\
\emph{facere} & \textbf{fazer} (to do) & hacer \\
\emph{ferrum} & \textbf{ferro} (iron) & hierro \\
\bottomrule
\end{tabularx}

So when a Spanish word starts with a silent \emph{h}, its Portuguese twin usually still has the honest Latin \emph{f} — \emph{hijo/filho}, \emph{hacer/fazer}. Portuguese is the more conservative sister here.
\end{cousinweb}

\begin{grammarlens}[title={Grammar Lens: the -ar present tense (pro-drop)}]
Drop \textbf{-ar} to get \emph{fal-}, then add the endings:

\noindent
\begin{tabularx}{\linewidth}{@{}>{\raggedright\arraybackslash}X>{\raggedright\arraybackslash}X>{\raggedright\arraybackslash}X@{}}
\toprule
\textbf{person} & \textbf{ending} & \textbf{\emph{falar} $\to$} \\
\midrule
(eu) & \textbf{-o} & \textbf{falo} (I speak) \\
(tu) & \textbf{-as} & \textbf{falas} (you speak) \\
(você/ele/ela) & \textbf{-a} & \textbf{fala} \\
(nós) & -amos & falamos \\
(vocês/eles) & -am & falam \\
\bottomrule
\end{tabularx}

Like Spanish and Italian, \textbf{Portuguese drops the subject pronoun} — \emph{falo} already says \textquotedblleft{}I,\textquotedblright{} so no \emph{eu} needed. (Portuguese joins the \emph{drop} camp: ES, PT, IT drop; FR, DE keep.)
\end{grammarlens}

\subsection*{Your turn}
Say these aloud:

\begin{itemize}
\raggedright
  \item \textquotedblleft{}falar\textquotedblright{} — \emph{fah-LAR}
  \item \textquotedblleft{}falo, falas, fala\textquotedblright{} — no \emph{eu} needed
  \item falar $\leftrightarrow$ hablar, filho $\leftrightarrow$ hijo — Portuguese keeps the \emph{f}, Spanish drops it to silent \emph{h}
\end{itemize}

\begin{tcolorbox}[breakable,colback=teal!4,colframe=teal!35!black,title={Before you move on}]
What Latin verb is \emph{falar} from, and which Spanish verb is its twin? (\emph{fabulārī} — Spanish \emph{hablar}.) What sound-law separates them? (Portuguese kept \emph{f-}; Spanish shifted \emph{f- $\to$ h-}: \emph{filho/hijo}.) Does Portuguese keep or drop the pronoun? (Drops it.) Next: \textbf{morar}, \textquotedblleft{}to live.\textquotedblright{}
\end{tcolorbox}
\section[morar]{morar — \textquotedblleft{}to live,\textquotedblright{} i.e. to \emph{linger} somewhere}

\label{lesson:PT-C05-morar}



\begin{quote}
A second \textbf{-ar} verb — and a distinctive Portuguese choice. Where its sisters say \textquotedblleft{}live\textquotedblright{} with \emph{habitāre} (\textquotedblleft{}keep having a place\textquotedblright{}), everyday Portuguese uses \textbf{morar} — from a verb that means to \textbf{linger}.
\end{quote}

\begin{sounds}
\begin{itemize}
\raggedright
  \item \textbf{morar} = \emph{moh-RAR}, tapped final \emph{r}, clean vowels.
\end{itemize}
\end{sounds}

\begin{cousinweb}
\textbf{morar} comes from Latin \textbf{morārī}, \emph{\textquotedblleft{}to delay, to linger, to tarry, to stay a while.\textquotedblright{}} To \emph{dwell} somewhere is, in Portuguese, to \emph{linger} there — a softer, more temporary image than \textquotedblleft{}keeping\textquotedblright{} a place. That \emph{morārī} root survives in English in a couple of unexpected words:

\begin{itemize}
\raggedright
  \item \textbf{moratorium} — an official \emph{delay} (a pause on payments or actions).
  \item \textbf{demur} — to \emph{linger / hesitate}, to raise an objection (\emph{without demur}).
\end{itemize}

(Portuguese does also have \emph{habitar} — the \emph{habitāre} verb its sisters use — but it's formal; \textbf{morar} is what you actually say: \emph{Onde você mora?}, \textquotedblleft{}Where do you live?\textquotedblright{})
\end{cousinweb}

\begin{grammarlens}[title={Grammar Lens: same -ar template}]
| (eu) moro · (tu) moras · (você) mora | (-o / -as / -a) |

| moramos · moram |  |

\begin{quote}
\textbf{Moro em Lisboa.} — \textquotedblleft{}I live in Lisbon.\textquotedblright{}
\end{quote}
\end{grammarlens}

\subsection*{Your turn}
Say these aloud:

\begin{itemize}
\raggedright
  \item \textquotedblleft{}morar\textquotedblright{} — \emph{moh-RAR}
  \item \textquotedblleft{}moro, moras, mora\textquotedblright{}
  \item \textquotedblleft{}Onde você mora?\textquotedblright{} — \textquotedblleft{}Where do you live?\textquotedblright{}; \textquotedblleft{}morar\textquotedblright{} $\leftrightarrow$ \emph{moratorium}
\end{itemize}

\begin{tcolorbox}[breakable,colback=teal!4,colframe=teal!35!black,title={Before you move on}]
What Latin verb is \emph{morar} from, and what did it mean? (\emph{morārī} — \textquotedblleft{}to delay / linger.\textquotedblright{}) English cousins? (\emph{Moratorium}, \emph{demur}.) Say \textquotedblleft{}I live in Lisbon.\textquotedblright{} (\emph{Moro em Lisboa.}) Next: \textbf{trabalhar}, \textquotedblleft{}to work.\textquotedblright{}
\end{tcolorbox}
\section[trabalhar]{trabalhar — \textquotedblleft{}to work,\textquotedblright{} and the torture returns}

\label{lesson:PT-C05-trabalhar}



\begin{quote}
A third \textbf{-ar} verb — and Portuguese rejoins its Iberian sister on one of the darkest etymologies anywhere. \textbf{trabalhar} (\textquotedblleft{}to work\textquotedblright{}), like Spanish \emph{trabajar}, began as a word for \textbf{torture}.
\end{quote}

\begin{sounds}
\begin{itemize}
\raggedright
  \item \texttt{nasal-anha} — the \textbf{-lh-} is a \emph{ly} glide (as in Spanish \emph{ll}): \textbf{trabalhar} = \emph{tra-ba-\textbf{LYAR}}.
\end{itemize}
\end{sounds}

\begin{cousinweb}
\textbf{trabalhar} comes from Vulgar Latin \textbf{tripaliāre}, \emph{\textquotedblleft{}to torture\textquotedblright{}} — from \textbf{tripalium}, a Roman instrument of torture made of \textbf{three stakes} (\emph{tri-} \textquotedblleft{}three\textquotedblright{} + \emph{pālus} \textquotedblleft{}stake\textquotedblright{} $\to$ English \emph{pale}, \emph{impale}). \emph{Torture $\to$ toil $\to$ work}, and English took the same road: \textbf{travail} (painful effort) $\to$ \textbf{travel} (a journey was an ordeal).

So Portuguese and Spanish are \textbf{twins} here — \emph{trabalhar} and \emph{trabajar}, both from \emph{tripaliāre}. Only the middle consonant differs: Latin \emph{-li-} became Portuguese \textbf{-lh-} (\emph{trabalhar}) and Spanish \textbf{-j-} (\emph{trabajar}). And this puts the \textquotedblleft{}work\textquotedblright{} verbs of the five languages into two neat camps:

\begin{itemize}
\raggedright
  \item \textbf{Torture} (\emph{tripalium}): Portuguese \emph{trabalhar}, Spanish \emph{trabajar}, French \emph{travailler}.
  \item \textbf{Labour} (\emph{labor}): Italian \emph{lavorare}.
\end{itemize}

Three languages remember work as \emph{torture}; only Italy remembers it as \emph{labour}.
\end{cousinweb}

\begin{grammarlens}[title={Grammar Lens: same -ar template}]
| (eu) trabalho · (tu) trabalhas · (você) trabalha | (-o / -as / -a) |

\begin{quote}
\textbf{Trabalho no Porto.} — \textquotedblleft{}I work in Porto.\textquotedblright{}
\end{quote}
\end{grammarlens}

\subsection*{Your turn}
Say these aloud:

\begin{itemize}
\raggedright
  \item \textquotedblleft{}trabalhar\textquotedblright{} — \emph{tra-ba-LYAR}
  \item \textquotedblleft{}trabalho, trabalhas, trabalha\textquotedblright{}
  \item trabalhar $\leftrightarrow$ trabajar (twins) $\to$ travail $\to$ travel — torture to journey
\end{itemize}

\begin{tcolorbox}[breakable,colback=teal!4,colframe=teal!35!black,title={Before you move on}]
What did \emph{trabalhar} originally mean, and from what device? (\textquotedblleft{}To torture,\textquotedblright{} $\leftarrow$ \emph{tripalium}, three stakes.) Its Spanish twin, and the one difference? (\emph{trabajar}; Latin \emph{-li-} $\to$ PT \emph{-lh-} / ES \emph{-j-}.) Which language uses \textquotedblleft{}labour\textquotedblright{} instead? (Italian, \emph{lavorare}.) Next: your first sentence — \textbf{Falo português}.
\end{tcolorbox}
\section[Falo português]{Falo português — your first real sentence}

\label{lesson:PT-C05-falo-portugues}



\begin{quote}
The payoff: you take a \textbf{verb} you conjugated and a \textbf{noun}, and build a sentence no one handed you — \emph{Falo português}, \textquotedblleft{}I speak Portuguese.\textquotedblright{} (\emph{Falo} covers both \textquotedblleft{}I speak\textquotedblright{} and \textquotedblleft{}I'm speaking.\textquotedblright{})
\end{quote}

\subsection*{You'll want to know: The new word: português}
\textbf{português} = \textquotedblleft{}Portuguese,\textquotedblright{} from \textbf{Portus Cale} — a Roman-era port at the mouth of the \textbf{Douro} river (today's \textbf{Porto} / \textbf{Oporto}). \emph{Portus} is Latin \textquotedblleft{}port, harbour\textquotedblright{} ($\to$ English \textbf{port}, \textbf{portal}, \textbf{import}); \textbf{Cale} is older and debated — perhaps Celtic for \textquotedblleft{}harbour/shelter,\textquotedblright{} or Latin \emph{cālidus} \textquotedblleft{}warm.\textquotedblright{} So \emph{Portugal} began as a single \textbf{harbour town}, whose name spread to the whole country and its language. Say it \emph{por-too-\textbf{GESH}} (final \emph{-s} = \emph{sh} in Portugal).

\subsection*{You'll want to know: The sentence, assembled}
\begin{quote}
\textbf{Falo} (I speak, from \emph{falar}) + \textbf{português} = \textbf{Falo português.}
\end{quote}

A complete, grammatical Portuguese sentence. As before:

\begin{itemize}
\raggedright
  \item \textbf{No \emph{eu}.} The \emph{-o} of \emph{falo} says \textquotedblleft{}I\textquotedblright{} — Portuguese drops the pronoun (like Spanish and Italian).
  \item \textbf{No article} on the language: \emph{falo português}.
\end{itemize}

With this, \textbf{all five languages} you've been building now cross the same line — from reciting phrases to \textbf{making sentences}:

\noindent
\begin{tabularx}{\linewidth}{@{}>{\raggedright\arraybackslash}X>{\raggedright\arraybackslash}X@{}}
\toprule
\textbf{} & \textbf{\textquotedblleft{}I speak/learn {[}language{]}\textquotedblright{}} \\
\midrule
Spanish & \textbf{Hablo} español \\
Portuguese & \textbf{Falo} português \\
Italian & \textbf{Parlo} italiano \\
French & \textbf{Je parle} français \\
German & \textbf{Ich lerne} Deutsch \\
\bottomrule
\end{tabularx}

Swap the pieces: \textbf{Moro em Lisboa.} · \textbf{Trabalho no Porto.} · \textbf{Você fala português?}

\subsection*{Your turn}
Say these aloud:

\begin{itemize}
\raggedright
  \item \textquotedblleft{}Falo português\textquotedblright{} — \emph{FAH-loo por-too-GESH}
  \item drop the \emph{eu} — Portuguese doesn't need it
  \item \textquotedblleft{}português\textquotedblright{} $\leftarrow$ \emph{Portus Cale}, a single harbour town that named a nation
\end{itemize}

\begin{tcolorbox}[breakable,colback=teal!4,colframe=teal!35!black,title={Before you move on}]
Where does \emph{português} come from? (\emph{Portus Cale}, a port at the Douro's mouth $\to$ Porto.) In \emph{Falo português}, why no \emph{eu} and no article? (The \emph{-o} is \textquotedblleft{}I\textquotedblright{}; languages take no article.) Which three of the five languages drop the pronoun? (Spanish, Portuguese, Italian.) Next: practice.
\end{tcolorbox}
\section[Practice]{Practice — making sentences with -ar verbs}

\label{lesson:PT-C05-practice}



\begin{quote}
For the first time you're \textbf{building} Portuguese sentences from a pattern. Drill the three verbs until the endings are automatic — and remember, Portuguese drops the \emph{eu}.
\end{quote}

\subsection*{Your turn: conjugate on command}
Run each through \emph{eu / tu / você} (pronoun dropped in speech):

\begin{itemize}
\raggedright
  \item \textbf{falar}: falo · falas · fala
  \item \textbf{morar}: moro · moras · mora
  \item \textbf{trabalhar}: trabalho · trabalhas · trabalha
\end{itemize}

Same endings every time (\textbf{-o / -as / -a}), no pronoun needed.

\subsection*{Your turn: say something real}
\begin{quote}
— \textbf{Você fala} português? — \textbf{Falo} um pouco. E você? — Também. \textbf{Onde} você \textbf{mora}? — \textbf{Moro} em Lisboa, e \textbf{trabalho} no Porto. — Obrigado! — De nada.
\end{quote}

\begin{grammarlens}[title={What you've built this chapter}]
\begin{itemize}
\raggedright
  \item \textbf{the regular -ar present tense} — drop \emph{-ar}, add \emph{-o/-as/-a/-amos/-am}; pronoun \textbf{dropped} (like Spanish, Italian).
  \item \textbf{falar} ($\leftarrow$ \emph{fabulārī} $\to$ fable; twin of Spanish \emph{hablar} but with the \emph{f} kept), \textbf{morar} ($\leftarrow$ \emph{morārī} \textquotedblleft{}to linger\textquotedblright{} $\to$ moratorium — a root no sibling uses), \textbf{trabalhar} ($\leftarrow$ \emph{tripalium} \textquotedblleft{}torture\textquotedblright{}; twin of \emph{trabajar}).
  \item \textbf{Falo português} ($\leftarrow$ \emph{Portus Cale}, a harbour town) — first self-assembled sentence, and the \textbf{fifth language} to cross from phrases into sentences.
\end{itemize}
\end{grammarlens}

\subsection*{Your turn}
Say these aloud:

\begin{itemize}
\raggedright
  \item conjugate all three verbs, eu/tu/você — no pronoun spoken
  \item answer about yourself — what you speak, where you live, where you work
  \item chain it — \textquotedblleft{}Olá! Falo português. Você fala português? … Até logo!\textquotedblright{}
\end{itemize}

\emph{Twice through:} \textquotedblleft{}Você fala português? — Falo um pouco.\textquotedblright{}

\begin{tcolorbox}[breakable,colback=teal!4,colframe=teal!35!black,title={Before you move on}]
What are the three singular \emph{-ar} endings? (\emph{-o / -as / -a}.) Give the \textquotedblleft{}I\textquotedblright{} forms of the three verbs. (\emph{Falo, moro, trabalho}.) \emph{falar} vs Spanish \emph{hablar} — what's the sound-law difference? (Portuguese kept \emph{f-}; Spanish shifted \emph{f$\to$h}.) Portuguese now moves in real sentences — this book has handed over rules, not phrases.
\end{tcolorbox}
